% ch4.tex -- Module 5 ESSENTIALS, Chapter 4: Gravity, orbits and Kepler's laws
% Condensed from src/mod5/ch4.tex (knowledge points only, no questions).
% Dot points: M5.IQ3.01, M5.IQ3.02, M5.IQ3.03, M5.IQ3.04
% Numbers: the full book's verified values; re-checked by
% Figures: figures/ch4_cannonball.pdf, ch4_equalareas.pdf, ch4_ellipseforces.pdf
%          (copied from mod5/figures with their scripts)
\chapter{Gravity, Orbits and Kepler's Laws}\label{ch4:chap}
\covers{M5.IQ3.01, M5.IQ3.02, M5.IQ3.03, M5.IQ3.04}

\begin{summarybox}{In this chapter}
\begin{itemize}
  \item Newton's law of universal gravitation and the field strength $g$.
  \item Why a satellite stays in orbit, and its speed, period and radius.
  \item Kepler's three laws, explained with forces and energy.
  \item Low Earth orbits and geostationary orbits, and their uses.
\end{itemize}
\end{summarybox}

%% ==================================================================
\section{Newton's law of gravitation and field strength}\label{ch4:sec:newton}
\covers{M5.IQ3.01}

\begin{bigidea}
All masses attract each other. The pull weakens quickly with distance.
\end{bigidea}

\begin{keywords}
  \kw{Gravitational force}{the pull between any two masses. Unit: newton (N).}
  \kw{Inverse-square law}{double the distance, and the effect is four times smaller.}
  \kw{Field strength ($g$)}{the gravitational force on each kilogram at a point.}
\end{keywords}

\begin{formulacard}{Newton's law of universal gravitation}
  \cardformula{F = \frac{GMm}{r^2}}
  \begin{symbols}
    \sym{$F$}{force on each object (equal and opposite)}{\si{N}}
    \sym{$G$}{universal gravitational constant, \num{6.67e-11}}{\si{N.m^2.kg^{-2}}}
    \sym{$M$, $m$}{masses of the two objects}{\si{kg}}
    \sym{$r$}{distance between the two centres of mass}{\si{m}}
  \end{symbols}
  \whentouse{any two masses (the law is \term{universal}: it works everywhere), treated as
    points or spheres, such as planets and moons.}
  \trap{$r$ is centre to centre, not surface to surface. Square all of $r$.}
\end{formulacard}

\subsection{Gravitational field strength}
\begin{itemize}
  \item A \term{gravitational field} is a region where any mass feels a gravitational force.
        $g = F/m$, in \si{N.kg^{-1}} = \si{m.s^{-2}}. In \term{free fall} (gravity is the only
        force), $g$ is the acceleration.
  \item A planet's field is \term{radial}: field lines point to the centre and spread out, so the
        field weakens with distance. Near the surface it is nearly \term{uniform} (the same
        size and direction everywhere), $g = \qty{9.8}{m.s^{-2}}$.
  \item The \term{altitude} $h$ is the height above the surface. Gravity formulas need
        $r = r_\mathrm{E} + h$ from the centre (a definition, not on the sheet). Data sheet:
        $r_\mathrm{E} = \qty{6.371e6}{m}$, $M_\mathrm{E} = \qty{6.0e24}{kg}$.
\end{itemize}

\begin{formulacard}[notonsheet]{Gravitational field strength}
  \cardformula{g = \frac{GM}{r^2}}
  \begin{symbols}
    \sym{$g$}{field strength (free-fall acceleration)}{\si{N.kg^{-1}} = \si{m.s^{-2}}}
    \sym{$G$}{universal gravitational constant}{\si{N.m^2.kg^{-2}}}
    \sym{$M$}{mass of the planet (or moon, or star)}{\si{kg}}
    \sym{$r$}{distance from the planet's \emph{centre} to the point}{\si{m}}
  \end{symbols}
  \whentouse{at a planet's surface ($r$ = its radius) or above it ($r = r_\mathrm{E} + h$).}
  \derive{$mg = \dfrac{GMm}{r^2}$ (the same force, written two ways). Divide both sides by $m$.}
  \trap{never use the altitude $h$ alone as $r$. Change km to m.}
\end{formulacard}

\textbf{Factors that affect $g$.}
\begin{itemize}
  \item Only the planet's mass ($g \propto M$) and the distance from its centre ($g \propto 1/r^2$):
        $r \times 2 \Rightarrow g \div 4$. The falling object's mass cancels, so all objects fall
        with the same acceleration (if air resistance is negligible).
  \item Data-sheet values give $GM_\mathrm{E} = \qty{4.0e14}{m^3.s^{-2}}$ and a surface $g$ of
        \qty{9.86}{m.s^{-2}}. The sheet's 9.8 differs only because $M_\mathrm{E}$ is rounded.
  \item Real Earth: $g \approx \qty{9.78}{m.s^{-2}}$ at the equator and \qty{9.83}{m.s^{-2}} at
        the poles (the equator bulges, and the spin reduces the measured $g$ there). Mountains and dense rock
        change $g$ slightly.
  \item $g$ is a vector towards the mass. Add the fields of two bodies as vectors.
\end{itemize}

%% ==================================================================
\section{Satellites in circular orbits}\label{ch4:sec:orbits}
\covers{M5.IQ3.02, M5.IQ3.04}

\begin{bigidea}
A satellite is always falling, but it moves sideways so fast that the ground curves away
beneath it. Gravity provides the centripetal force.
\end{bigidea}

\begin{sidebyside}[0.46]
  \centering
  \includegraphics{ch4_cannonball.pdf}
  \captionof{figure}{Newton's cannonball (mountain far too tall). C is a circular orbit:
    gravity $F$ is at \ang{90} to $v$.}\label{ch4:fig:cannon}
\nextside
  Newton imagined a cannon on a very tall mountain, with no air (\cref{ch4:fig:cannon}).
  \begin{itemize}
    \item A slow ball falls to the ground (A). A faster one lands further away (B).
    \item At about \qty{7.9}{km.s^{-1}} near the surface, the ground curves away as fast as the
          ball falls. It never lands: it is in \term{orbit} (C).
    \item Faster still, the orbit is an oval (D). Above \qty{11.2}{km.s^{-1}} the ball escapes.
    \item An orbit is projectile motion that never ends: free fall. Gravity points to the
          centre, so \emph{gravity provides the centripetal force}: $F = F_\mathrm{c}$.
  \end{itemize}
\end{sidebyside}

\begin{formulacard}[notonsheet]{Orbital speed (circular orbit)}
  \cardformula{v = \sqrt{\frac{GM}{r}}}
  \begin{symbols}
    \sym{$v$}{orbital speed (the speed for a circular orbit)}{\si{m.s^{-1}}}
    \sym{$G$}{universal gravitational constant}{\si{N.m^2.kg^{-2}}}
    \sym{$M$}{mass of the central body (\emph{not} the satellite)}{\si{kg}}
    \sym{$r$}{orbital radius, from the centre of the central body}{\si{m}}
  \end{symbols}
  \whentouse{any circular orbit around a planet, moon or star.}
  \derive{$\dfrac{GMm}{r^2} = \dfrac{mv^2}{r}$ (both on the sheet). Divide by $m$, multiply by
    $r$: $v^2 = \dfrac{GM}{r}$. Take the square root.}
  \trap{$M$ is the planet's mass, not the satellite's. $r$ is not the altitude.}
\end{formulacard}

\begin{itemize}
  \item The satellite's mass cancels: all objects at one height orbit at the same speed, with the
        same period (the orbit simulation in \cref{ch4:sec:experiments} shows this).
  \item A higher orbit is slower: $v \propto 1/\sqrt r$. The acceleration
        $a_\mathrm{c} = v^2/r = GM/r^2$ is the local $g$.
  \item Orbital velocity is along the tangent to the orbit, at \ang{90} to the radius.
  \item \term{Apparent weightlessness} (feeling no weight in free fall): on the ISS (\qty{420}{km}
        up) $g = \qty{8.7}{m.s^{-2}}$, but astronauts and station fall together, so the floor does
        not push: the \term{normal force} is zero.
\end{itemize}

\subsection{The period and Kepler's third law}
\begin{formulacard}{Kepler's third law (Newton's form)}
  \cardformula{\frac{r^3}{T^2} = \frac{GM}{4\pi^2}}
  \begin{symbols}
    \sym{$r$}{orbital radius (for an ellipse: the semi-major axis)}{\si{m}}
    \sym{$T$}{period: the time for one orbit}{\si{s}}
    \sym{$G$}{universal gravitational constant}{\si{N.m^2.kg^{-2}}}
    \sym{$M$}{mass of the central body}{\si{kg}}
  \end{symbols}
  \whentouse{to find $T$, $r$ (then $h = r - r_\mathrm{E}$) or the central mass $M$.}
  \derive{put $v = \dfrac{2\pi r}{T}$ (on the sheet) into $v^2 = \dfrac{GM}{r}$:
    $\dfrac{4\pi^2 r^2}{T^2} = \dfrac{GM}{r}$. Multiply by $r$, divide by $4\pi^2$.}
  \trap{$T$ in seconds. Use a cube root (not a square root) to find $r$.}
\end{formulacard}

%% ==================================================================
\section{Kepler's laws of planetary motion}\label{ch4:sec:kepler}
\covers{M5.IQ3.04}

\begin{bigidea}
Kepler's laws came from data. Newton's law of gravitation explains all three.
\end{bigidea}

Kepler used Tycho Brahe's careful data. He published his first two laws in 1609 and the third
in 1619. They are \term{empirical laws}: they describe but do not explain. Newton (1687) showed that a force with
$F \propto 1/r^2$ explains all three.

\textbf{First law (ellipses).} \emph{Each planet moves in an ellipse, with the Sun at one focus.}
An \term{ellipse} is an oval with two \term{foci} (fixed points on its long axis). The other focus
is empty. Planet orbits are nearly circles (the Earth's \term{eccentricity}, its stretch,
is only 0.017), so they are treated as circles. Comets' orbits are very stretched. The closest and farthest points from
the Sun are \term{perihelion} and \term{aphelion} (around the Earth: perigee, apogee). For an
ellipse, $r$ in Kepler's third law is the \term{semi-major axis} (half the longest diameter),
$a = (r_\mathrm{p} + r_\mathrm{a})/2$.

\begin{sidebyside}[0.56]
  \centering
  \includegraphics{ch4_equalareas.pdf}
  \captionof{figure}{Equal times, equal areas: $\Delta t_1 = \Delta t_2$, so $A_1 = A_2$ (to
    scale).}\label{ch4:fig:areas}
\nextside
  \textbf{Second law (equal areas).} \emph{A line from the Sun to a planet sweeps out equal
  areas in equal times} (\cref{ch4:fig:areas}). Near the Sun the line is short, so the planet
  covers a longer arc in the same time: it is \textbf{fastest at perihelion},
  \textbf{slowest at aphelion}.
\end{sidebyside}

\begin{sidebyside}[0.52]
  \centering
  \includegraphics{ch4_ellipseforces.pdf}
  \captionof{figure}{Velocity (blue): length shows the speed. Gravity (red): direction only.
    Dashed red: the part of gravity along the motion.}\label{ch4:fig:forces}
\nextside
  \textbf{With forces} (\cref{ch4:fig:forces}). A force does \term{work} (gives energy) only if
  part of it points along the motion.
  \begin{itemize}
    \item \textbf{Circle:} gravity is at \ang{90} to $v$. No work, so constant speed.
    \item \textbf{Moving closer:} part of gravity points along the motion, so the planet
          speeds up.
    \item \textbf{Moving away:} part of gravity points against the motion, so it slows down.
  \end{itemize}
\end{sidebyside}

\textbf{With energy.} Only gravity acts, so the \term{total energy} (kinetic $K$ plus potential
$U$) is constant. $U = -\dfrac{GMm}{r}$ (on the sheet) is lowest (most negative) at perihelion, so
$K = \tfrac12 mv^2$ is largest there. In a circular orbit $U$, $K$ and the speed are constant.
At perihelion and aphelion only, $r_\mathrm{p} v_\mathrm{p} = r_\mathrm{a} v_\mathrm{a}$ (from
equal areas; not on the sheet; distances from the Sun).

\textbf{Third law (periods).} \emph{The square of the period is proportional to the cube of the
orbital radius (semi-major axis)}: $T^2 \propto r^3$. The constant, $GM/4\pi^2$, depends only on
the central mass: all the Sun's planets share one value, and all the Earth's satellites share
another. For two orbits of one body, $r_1^3/T_1^2 = r_2^3/T_2^2$ in any units, such as
years and \term{astronomical units} (1~AU = average Earth--Sun distance, \qty{1.50e11}{m}).

%% ==================================================================
\section{Low Earth orbits and geostationary orbits}\label{ch4:sec:leogeo}
\covers{M5.IQ3.03}

\begin{bigidea}
The height of an orbit fixes its speed and period. Each height suits different jobs
(\cref{ch4:tab:leogeo}).
\end{bigidea}

\begin{itemize}
  \item A \term{low Earth orbit} (LEO) is \qtyrange{200}{2000}{km} up. Thin air causes
        \term{orbital decay} (slow loss of height). A \emph{polar} LEO passes over both poles.
        The Earth turns beneath it, so it passes over every part of the Earth.
  \item A \term{geostationary} satellite stays above one point on the equator. It needs:
        (1) a period of one day (strictly 23\,h\,56\,min); (2) a circular orbit above the
        equator; (3) motion west to east, like the Earth's spin. Kepler's third law then allows
        only one radius, \qty{4.23e7}{m} (6.6 Earth radii).
  \item A \term{geosynchronous} orbit has a one-day period but may be tilted. Then the satellite
        drifts north and south in the sky, so it is not geostationary.
\end{itemize}

\begin{table}
  \caption{Low Earth orbit and geostationary orbit compared.}\label{ch4:tab:leogeo}
  \centering
  \small
  \begin{tabular}{@{}L{25mm}L{57mm}L{58mm}@{}}
    \toprule
    \colhead{Property} & \colhead{Low Earth orbit} & \colhead{Geostationary orbit}\\
    \midrule
    Altitude & \qtyrange{200}{2000}{km} (ISS: \qty{420}{km}) & about \qty{36000}{km}\\
    Speed, period & about \qty{7.7}{km.s^{-1}}, about \qty{90}{min} & \qty{3.1}{km.s^{-1}}, 1 day\\
    In the sky & crosses the sky in minutes & fixed, so ground dishes never move\\
    View & small area, fine detail & about 40\,\% of the Earth; poor near the poles\\
    Signal delay & a few milliseconds & about \qty{0.24}{s} (up and down)\\
    Problems & orbital decay; needs many satellites & weaker signals; costly launch\\
    Uses & ISS, imaging, mapping, polar weather satellites, internet constellations
      & TV, communications, weather images of a whole hemisphere\\
    \bottomrule
  \end{tabular}
\end{table}

Each use comes from a property: a low orbit gives fine detail \emph{because} it is close to the
ground. In both orbits gravity provides the centripetal force.

%% ==================================================================
% layout QA (after the whole-book build): without this break the first box started at the foot
% of the page with only its Aim and the "Method" heading. Section 4.5 (three boxes) plus the
% Remember box fill exactly two pages.
\clearpage
\section{Key experiments}\label{ch4:sec:experiments}
\covers{M5.IQ3.01, M5.IQ3.02, M5.IQ3.04}

\begin{experiment}{Factors that affect the field strength \texorpdfstring{$g$}{g}}
\aim{To show that surface $g$ is proportional to $M/r^2$ ($M$: mass, $r$: radius of the body).}
\method{\term{Secondary data} (measured by others) for the Moon, Mercury, Mars, Venus and the
  Earth; $g$ plotted against $M/r^2$.}
\results{A straight line through the origin; gradient \qty{6.67e-11}{N.m^2.kg^{-2}}.}
\conclusion{$g \propto M/r^2$, and the gradient is $G$, as $g = GM/r^2$ predicts.}
\significance{$M = gr^2/G$ gives a planet's mass. Cavendish found $G$ in 1798 (``weighing the Earth'').}
\end{experiment}

\begin{experiment}{Kepler's third law and the moons of Jupiter}
\aim{To show that all of Jupiter's large moons share one value of $r^3/T^2$, and to find Jupiter's mass.}
\method{$r$ and $T$ (in seconds) of Io, Europa, Ganymede and Callisto; $r^3$ plotted against $T^2$.}
\results{A straight line through the origin; gradient \qty{3.205e15}{m^3.s^{-2}}.}
\conclusion{$r^3 \propto T^2$. Gradient $= GM/4\pi^2$, so Jupiter's mass is \qty{1.90e27}{kg}.}
\significance{Galileo (1610): the first objects seen orbiting another planet. This supported the
  Sun-centred model.}
\end{experiment}

% Orbit simulation: the full book's numbers (mod5/ch4.tex, "Orbits in a simulation");
\begin{experiment}{Orbits in a simulation: mass, radius, speed and period}
\aim{To find how a satellite's period and speed depend on its own mass and on its orbital
  radius.}
\method{A computer orbit model (planet mass \qty{6.0e24}{kg}). The satellite is launched at
  \ang{90} to the radius, at the speed that gives a circular path; 5 orbits are timed. Only the
  satellite's mass is changed (\qty{100}{kg}, \qty{1000}{kg}, \qty{5000}{kg}), then only $r$
  (\num{7.0e6} to \qty{4.0e7}{m}).}
\results{The same period for every mass ($T = \qty{5817}{s}$ at $r = \qty{7.0e6}{m}$). A bigger
  $r$ gave a longer $T$ and a lower $v$, but $v^2 r = \qty{4.00e14}{m^3.s^{-2}}$ at every radius.}
\conclusion{The period does not depend on the satellite's mass. $v^2 r = GM$, so
  $v \propto 1/\sqrt{r}$, as $v = \sqrt{GM/r}$ predicts.}
\significance{All objects at one height move together, whatever their mass: an astronaut floats
  beside the space station.}
\end{experiment}

%% ==================================================================
\begin{summarybox}{Remember}
\begin{itemize}
  \item $F = GMm/r^2$ (on the sheet), $g = GM/r^2$ (not on it). $r = r_\mathrm{E} + h$, in metres.
  \item Gravity provides the centripetal force: $v = \sqrt{GM/r}$ (not on the sheet).
  \item A higher orbit is slower. The satellite's mass cancels.
  \item Weightless astronauts fall with their craft: the normal force is zero, not gravity.
  \item $r^3/T^2 = GM/(4\pi^2)$ for all orbits around one body, with $T$ in seconds.
  \item Kepler: ellipse with the Sun at a focus; equal areas, so fastest at perihelion; $T^2 \propto r^3$.
  \item Elliptical orbit: gravity does work, but $K + U$ stays constant ($K$ largest at perihelion).
  \item LEO: low, about \qty{90}{min}, fine detail. Geostationary: \qty{36000}{km}, one day, above
        the equator, west to east.
\end{itemize}
\end{summarybox}
